% from ega0-21-fr.tex
\oldpage[0\textsubscript{IV-1}]{178}
en vertu de (21.9.3) $\widehat\Omega^1_{A_0/k}$ est un $A_0$-module libre de
rang $r=\dim(A_0)=\dim(A)$ (16.1.5)~; le produit tensoriel
$\widehat\Omega^1_{A_0/k}\otimes_{A_0}A$ est donc un $A$-module libre de rang
$r$ identique à son séparé complété~; enfin $\widehat\Omega^1_{A/A_0}$ est
un $A$-module de type fini (20.4.7), donc identique à son séparé complété
($\mathbf 0_{\mathrm I}$, 7.3.6). Cela étant, dans la suite d'homomorphismes
\begin{equation}
\tag{21.9.5.2}
\label{0.21.9.5.2-fr}
  \widehat\Omega^1_{A_0/k}\otimes_{A_0}A
  \xrightarrow{u}\widehat\Omega^1_{A/k}
  \xrightarrow{v}\widehat\Omega^1_{A/A_0}\longrightarrow0
\end{equation}
on sait que $v$ est surjectif et que $\operatorname{Im}(u)$ est dense dans
$\operatorname{Ker}(v)$ (20.7.17)~; mais tout sous-$A_0$-module de
$\widehat\Omega^1_{A/k}$ est fermé pour la topologie $\mathfrak m_0$-adique
($\mathbf 0_{\mathrm I}$, 7.3.5), donc la suite (21.9.5.2) est exacte. Notons
d'autre part que puisque $A$ est entier sur $A_0$, $A_0$ intégralement clos
et $E$ séparable sur $L_0$, $\widehat\Omega^1_{A/A_0}$ est un $A$-module de
torsion (20.4.13, (iv))~; tensorisant la suite exacte (21.9.5.2) par $E$, il
vient une suite exacte
\[
  \widehat\Omega^1_{A_0/k}\otimes_{A_0}E
  \longrightarrow \widehat\Omega^1_{A/k}\otimes_AE\longrightarrow0
\]
d'où $\operatorname{rg}_E(\widehat\Omega^1_{A/k}\otimes_AE)\leqslant r$. Il
reste donc à montrer qu'il existe, dans $\widehat\Omega^1_{A/k}\otimes_AE$,
$r$ éléments linéairement indépendants sur $E$. Or, soient
$D_i=\partial/\partial T_i$ les dérivations canoniques de $A_0$ dans lui-même
$(1\leqslant i\leqslant r)$~; elles se prolongent de façon unique en des
dérivations (encore notées $D_i$) de $E$ dans lui-même, puisque $E$ est
extension séparable finie de $L_0$~; par restriction à $A$, ces dérivations
donnent des dérivations de $A$ dans $E$, et il est immédiat que ces
dérivations prennent leurs valeurs dans un même sous-$A$-module de type fini
de $E$~; comme elles sont continues dans $A_0$ et que la topologie de $A$ est
la topologie $\mathfrak m_0$-adique, on a ainsi formé $r$ dérivations
continues de $A$ dans un $A$-module topologique, qui sont évidemment
linéairement indépendantes puisque $D_i(T_j)=\delta_{ij}$~; ceci achève de
prouver la proposition.
\end{proof}

La proposition suivante est un analogue «~formel~» de (21.7.3)~:

\begin{proposition}[21.9.6]
\label{0.21.9.6-fr}
Soient $k_0\subset k\subset K_0$ trois corps de caractéristique $p>0$, tels
que $[K_0:k]<+\infty$~; on pose
\[
  L_0=K_0((T_1,\ldots,T_r)),\qquad
  M_0=K_0((T_1^p,\ldots,T_r^p)),
\]
\[
  L=k((T_1,\ldots,T_r)),\qquad
  M=k((T_1^p,\ldots,T_r^p)),\qquad
  N=k((T_1^{p^2},\ldots,T_r^{p^2})).
\]
Soit $E$ une extension de type fini de $L_0$.
\begin{enumerate}
  \item[(i)] Si $\Upsilon_{K_0/k_0}$ est de rang fini, on a
  \begin{equation}
  \tag{21.9.6.1}
  \label{0.21.9.6.1-fr}
    \operatorname{rg}_E\Omega^1_{E/M}
    -\operatorname{rg}_{K_0}\Omega^1_{K_0/k}
    =(r+\deg.\operatorname{tr}_{L_0}E)
      +d(E/K_0,k_0)+d(E/M_0,N/k_0)
  \end{equation}
  (notations de (21.6.1)), et en particulier on a
  \begin{equation}
  \tag{21.9.6.2}
  \label{0.21.9.6.2-fr}
    \operatorname{rg}_E\Omega^1_{E/M}
    -\operatorname{rg}_{K_0}\Omega^1_{K_0/k}
    \geqslant r+\deg.\operatorname{tr}_{L_0}E+d(E/K_0,k_0).
  \end{equation}
  En outre, si les deux membres de (21.9.6.2) sont égaux, ils le restent
  lorsqu'on remplace $k$ par un corps $k'$ tel que $k_0\subset k'\subset k$.

  \item[(ii)] On suppose en outre que $[k_0:k_0^p]<+\infty$. Soit
  $(k_\alpha)$ une famille filtrante décroissante de sous-corps de $K_0$
  contenant $k_0$, tels que $[K_0:k_\alpha]<+\infty$ pour tout $\alpha$ et
  que
  $\bigcap_\alpha k_\alpha(K_0^p)=k_0(K_0^p)$~; alors il existe un $\alpha$
  tel que, pour $k=k_\alpha$, les deux membres de (21.9.6.2) soient égaux.
\end{enumerate}
\end{proposition}

\begin{proof}
\begin{enumerate}
  